\section{Заключение}
\label{sec:conclusion} \index{Conclusion}

В рамках данной работы были решены все поставленные задачи:

\begin{itemize}
\item Изучена генерация отладочной информации для архитектуры RISC-V компиляторами GCC и LLVM

В рамках этой задачи была подробная рассмотрена архитектура RISC-V, ее особенности и преимущества.
Также, было подробно изучено векторное расширение данной архитектуры и специфика соглашения о
вызовах для векторных типов.
\item Реализовано векторное соглашение о вызовах в отладчике GDB
\item Поддержано определение векторных типов в отладчике GDB
\end{itemize}

Была выполнена основная цель работы: реализовать в отладчике GDB для архитектуры RISC-V
поддержку вызова функций, использующих векторный ABI.

Корректность работы данной реализации была проверена на тестовых сценариях в рамках
тестовой инфраструктуры GDB.

Патч с поддержкой функциональности планируется к включению в Syntacore Development Toolkit (SC-DT)
и публикации в открытом доступе.

В рамках дальнейшей работы планируется расширить спектр поддерживаемых функциональностей в
отладчике GDB для векторного расширения RISC-V. Одним из направлений может стать
реализация обратного исполнения с учетом специфики векторного расширения RISC-V.

\newpage
